% R4 regime diagram: two panels, transcribed from figures/src/r4_figures.py
% (the second figure block: axA/axB). Panel A: R(delta,f) curves at
% delta in {0, 0.01, 0.04}, from the same rate_general closed form used
% throughout paper1 (recomputed here, 30 points/curve). Panel B: the zoom
% advantage as a function of flagged-set density d = k/F.
%
% PANEL B REPLOTTED (not reproduced): the script plots ONLY the idealized
% advantage F/(k log2(F/k)) = 1/(d*log2(1/d)). That quantity has a global
% minimum of 1.884x at d = 1/e and max_d d*log2(1/d) = 0.531 < 1, so it
% CANNOT cross parity anywhere on (0,1) -- extending the axis would not help.
% The panel therefore shaded a "zoom wins" band across 100% of the chart and
% drew a "zoom stops paying" line the plotted curve was mathematically
% incapable of touching, while paper1.tex Sec. sec:zoom spends a paragraph on
% the deployed algorithm LOSING at high density. The figure's graphical
% grammar asserted the opposite of the section's own boundary, in a paper
% whose limitations box commits to bounding "the deployed algorithm, not the
% best one".
%
% The primary series is now the DEPLOYED guaranteed advantage
%   adv(d) = 1/(d*(2*ceil(log2(1/d)) + 4)),
% Theorem 4's own bound rewritten in density (it is a function of d alone,
% and returns exactly 12.500x at d = 0.004, reproducing the Corollary as a
% built-in cross-check). It crosses parity at d = 1/12 = 0.08333, well inside
% the plotted x-range, so the shaded region is now clipped there instead of
% filling the chart, and the parity line finally has something that reaches
% it. The idealized curve is kept as a faint dashed reference -- the gap
% between the two IS the factor-~2 that sec:zoom's "Positioning against group
% testing" paragraph is about -- and extended to d = 0.95 so the paper's own
% (F,k) = (100,90) illustration (d = 0.9, 199 opens against 100 flat) is on
% the axes rather than off the right edge. Also marked: the two operating
% points Sec. sec:rdf actually recommends, which sit on the LOSING side.
% The spurious "F=2500" in the panel title is dropped -- the curve has no
% dependence on F beyond the measured point.
% All Panel B coordinates recomputed from the two closed forms. [verified]
%
% BUG FOUND AND CORRECTED (not silently): the script's Panel A calls the raw
% pinned formula `rate_general` across the WHOLE plotted f range for every
% delta, without switching to true_rate's R=0 correction past the
% Proposition's own boundary f=1-delta/p. For delta=0.04 that boundary is
% f*=0.2, well inside the plotted range [0.04,0.4] -- so the script's raw
% curve dips to 0 near f=0.19-0.21 and then rises again to 0.0069 by f=0.4,
% i.e. it reproduces EXACTLY the "spurious uptick past the domain of
% validity" bug that paper1.tex's own text (Sec. sec:rdf, and this figure's
% neighbor fig:rdffrontier) says was already caught and fixed. It was not
% fixed here. Per the house falsification-first rule this is corrected in
% the curve below (delta=0.04 is clipped to R=0 for f>=0.2, matching
% true_rate/paper1_rdf_frontier_figures.py's already-correct logic) rather
% than reproduced; the .py script itself is left untouched, and this is
% flagged prominently in the figure-conversion report.
% [verified: script r4_figures.py, correction applied per CONVENTION.md]
\begin{figure}[t]
\centering
\textbf{\color{harborblue}R4 regime --- the two-constraint dial and the boundary where zooming stops paying}\par\vspace{4pt}
\begin{minipage}{0.48\textwidth}
\centering
\begin{tikzpicture}
\begin{axis}[
  harbor curve,
  height=7.2cm,
  xlabel={flag rate $f$},
  ylabel={rate $R(\delta,f)$ (bits/sym)},
  title={Panel A --- the two-constraint rate surface\\(Bernoulli $p{=}0.05$)},
  xmin=0,xmax=0.4,ymin=0,ymax=0.34,
  legend pos=north east,
  legend style={font=\scriptsize,draw=black!25,fill=white,fill opacity=0.92,text opacity=1,rounded corners=1pt,inner sep=3pt,title={\scriptsize miss tolerance}},
]
\addplot[harborblue,thick] coordinates {
(0.0500,0.2864) (0.0621,0.2423) (0.0741,0.2189) (0.0862,0.2018) (0.0983,0.1881) (0.1103,0.1768) (0.1224,0.1670) (0.1345,0.1584) (0.1466,0.1507) (0.1586,0.1438) (0.1707,0.1375) (0.1828,0.1317) (0.1948,0.1263) (0.2069,0.1213) (0.2190,0.1167) (0.2310,0.1123) (0.2431,0.1082) (0.2552,0.1043) (0.2672,0.1006) (0.2793,0.0970) (0.2914,0.0937) (0.3034,0.0905) (0.3155,0.0874) (0.3276,0.0845) (0.3397,0.0816) (0.3517,0.0789) (0.3638,0.0763) (0.3759,0.0738) (0.3879,0.0713) (0.4000,0.0690)
};
\addlegendentry{$\delta=0.00$}
\addplot[seagreen,dashed,thick] coordinates {
(0.0400,0.2062) (0.0524,0.1650) (0.0648,0.1443) (0.0772,0.1296) (0.0897,0.1181) (0.1021,0.1086) (0.1145,0.1005) (0.1269,0.0935) (0.1393,0.0873) (0.1517,0.0817) (0.1641,0.0767) (0.1766,0.0721) (0.1890,0.0679) (0.2014,0.0640) (0.2138,0.0604) (0.2262,0.0571) (0.2386,0.0539) (0.2510,0.0510) (0.2634,0.0482) (0.2759,0.0455) (0.2883,0.0431) (0.3007,0.0407) (0.3131,0.0385) (0.3255,0.0363) (0.3379,0.0343) (0.3503,0.0324) (0.3628,0.0305) (0.3752,0.0287) (0.3876,0.0270) (0.4000,0.0254)
};
\addlegendentry{$\delta=0.01$}
\addplot[shipred,dashdotted,thick] coordinates {
(0.0400,0.0141) (0.0524,0.0104) (0.0648,0.0078) (0.0772,0.0059) (0.0897,0.0044) (0.1021,0.0032) (0.1145,0.0023) (0.1269,0.0016) (0.1393,0.0011) (0.1517,0.0006) (0.1641,0.0003) (0.1766,0.0001) (0.1890,0.0000) (0.2014,0.0000) (0.2138,0.0000) (0.2262,0.0000) (0.2386,0.0000) (0.2510,0.0000) (0.2634,0.0000) (0.2759,0.0000) (0.2883,0.0000) (0.3007,0.0000) (0.3131,0.0000) (0.3255,0.0000) (0.3379,0.0000) (0.3503,0.0000) (0.3628,0.0000) (0.3752,0.0000) (0.3876,0.0000) (0.4000,0.0000)
};
\addlegendentry{$\delta=0.04$}
\addplot[only marks,mark=diamond*,mark size=2.6pt,harborblue,mark options={draw=black,line width=0.5pt},forget plot] coordinates {(0.05,0.2864)};
\node[regimebox,fill=white,fill opacity=0.9,draw=harborblue,font=\scriptsize,inner sep=1.5pt,anchor=east] at (axis cs:0.39,0.20)
  {$H(0.05)=0.286$};
\draw[->,harborblue,thick] (axis cs:0.32,0.19) -- (axis cs:0.06,0.278);
\addplot[only marks,mark=*,mark size=2pt,harborblue,mark options={draw=black,line width=0.4pt},forget plot] coordinates {(0.10,0.1864)};
\node[font=\scriptsize,harborblue,fill=white,fill opacity=0.9,inner sep=1.5pt,anchor=east] at (axis cs:0.39,0.14) {$R(0.00,0.10)=0.186$};
\draw[->,harborblue] (axis cs:0.32,0.135) -- (axis cs:0.107,0.188);
\addplot[only marks,mark=*,mark size=2pt,seagreen,mark options={draw=black,line width=0.4pt},forget plot] coordinates {(0.10,0.1100)};
\node[font=\scriptsize,seagreen,fill=white,fill opacity=0.9,inner sep=1.5pt,anchor=east] at (axis cs:0.39,0.08) {$R(0.01,0.10)=0.110$};
\draw[->,seagreen] (axis cs:0.32,0.078) -- (axis cs:0.108,0.113);
\addplot[only marks,mark=*,mark size=2pt,shipred,mark options={draw=black,line width=0.4pt},forget plot] coordinates {(0.06,0.009)};
\node[font=\scriptsize,shipred,fill=white,fill opacity=0.9,inner sep=1.5pt,anchor=east] at (axis cs:0.39,0.02) {$R(0.04,0.06)=0.009$};
\draw[->,shipred] (axis cs:0.32,0.019) -- (axis cs:0.065,0.012);
\end{axis}
\end{tikzpicture}
\end{minipage}%
\hfill
\begin{minipage}{0.48\textwidth}
\centering
\begin{tikzpicture}
\begin{axis}[
  harbor curve,
  xmode=log,ymode=log,
  height=7.2cm,
  xlabel={flagged-set density $d=k/F$},
  ylabel={open-count advantage (log)},
  title={Panel B --- the boundary, deployed:\\zoom pays only for $d\le 1/12$},
  xmin=0.0008,xmax=1.0,ymin=0.10,ymax=130,
  legend pos=north east,
  legend style={font=\tiny,draw=black!25,fill=white,fill opacity=0.92,text opacity=1,rounded corners=1pt,inner sep=2pt},
]
% Shaded region: where the DEPLOYED bound actually exceeds parity, i.e.
% d <= 1/12. Follows the deployed curve down to (1/12, 1) and back along y=1.
\addplot[draw=none,fill=harborblue,fill opacity=0.12] coordinates {
(0.0010000,41.6667) (0.0019531,21.3333) (0.0019531,23.2727) (0.0039063,11.6364)
(0.0039063,12.8000) (0.0078125,6.4000) (0.0078125,7.1111) (0.0156250,3.5556)
(0.0156250,4.0000) (0.0312500,2.0000) (0.0312500,2.2857) (0.0625000,1.1429)
(0.0625000,1.3333) (0.0833333,1.0000)
(0.0833333,1) (0.0010000,1)
} \closedcycle;
\addlegendentry{zoom wins (deployed bound $>1$)}
% The deployed guaranteed advantage 1/(d(2*ceil(log2(1/d))+4)). Exact: within
% each octave [2^-j, 2^-(j-1)) the ceiling is constant, so the curve is a
% straight line of slope -1 on log-log and two endpoints per octave are
% exact; it steps UP at each octave boundary as the ceiling drops.
\addplot[harborblue,thick,no marks] coordinates {
(0.0010000,41.6667) (0.0019531,21.3333)
(0.0019531,23.2727) (0.0039063,11.6364)
(0.0039063,12.8000) (0.0078125,6.4000)
(0.0078125,7.1111) (0.0156250,3.5556)
(0.0156250,4.0000) (0.0312500,2.0000)
(0.0312500,2.2857) (0.0625000,1.1429)
(0.0625000,1.3333) (0.1250000,0.6667)
(0.1250000,0.8000) (0.2500000,0.4000)
(0.2500000,0.5000) (0.5000000,0.2500)
(0.5000000,0.3333) (0.9500000,0.1754)
};
\addlegendentry{deployed: $\le 2k\lceil\log_2(F/k)\rceil{+}4k$ (Thm.~4)}
\addplot[harborblue!40,densely dashed,thick,no marks] coordinates {
(0.001000,100.3433) (0.001314,79.5156) (0.001726,63.1178) (0.002268,50.1940) (0.002980,39.9968) (0.003915,31.9415) (0.005144,25.5706) (0.006758,20.5255) (0.008879,16.5251) (0.011666,13.3489) (0.015327,10.8239) (0.020137,8.8142) (0.026458,7.2128) (0.034761,5.9359) (0.045671,4.9175) (0.060005,4.1060) (0.078837,3.4610) (0.103580,2.9513) (0.136089,2.5538) (0.178801,2.2519) (0.234917,2.0370) (0.308646,1.9104) (0.405515,1.8938) (0.532785,2.0663) (0.700000,2.7762) (0.750000,3.2125) (0.850000,5.0180) (0.900000,7.3098) (0.950000,14.2246)
};
\addlegendentry{idealized $k\log_2(F/k)$ --- never crosses parity}
\addplot[seagreen,dashed,thick,no marks] coordinates {(0.0008,1) (1.0,1)};
\addlegendentry{parity: zoom stops paying}
\addplot[only marks,mark=triangle*,mark size=3.0pt,shipred,mark options={draw=red!30!black,line width=0.5pt}] coordinates {(0.004,15.3) (0.9,0.5025)};
\addlegendentry{measured: $(F,k){=}(2500,10)$ and $(100,90)$}
\addplot[only marks,mark=square*,mark size=2.2pt,seagreen,mark options={draw=black,line width=0.4pt}] coordinates {(0.5,0.3333) (0.4,0.3125)};
\addlegendentry{\S\ref{sec:rdf} at $f{=}2p$ and $(\delta,f){=}(0.01,0.10)$}
\addplot[only marks,mark=*,mark size=2.4pt,shipred,mark options={draw=black,line width=0.4pt},forget plot] coordinates {(0.004,12.5) (0.0833333,1.0) (0.9,0.1852)};
\node[font=\tiny,text=shipred,anchor=north east] at (axis cs:0.076,0.93) {$d{=}1/12$};
\node[regimebox,fill=shipred!8,draw=shipred,font=\tiny,text width=2.3cm,inner sep=2pt,anchor=west] at (axis cs:0.00090,0.45)
  {$(2500,10)$: bound\\$12.5\times$, measured\\$15.4\times$, ideal $31\times$};
\draw[->,shipred,thick] (axis cs:0.0042,0.72) -- (axis cs:0.0042,10.8);
\node[font=\tiny,text=shipred,anchor=east,align=right] at (axis cs:0.85,0.145)
  {$(100,90)$: $199$ opens\\against $100$ flat};
\draw[->,shipred] (axis cs:0.86,0.16) -- (axis cs:0.895,0.46);
\end{axis}
\end{tikzpicture}
\end{minipage}
\caption{Regime diagram for the zoom advantage: the idealized open-count advantage $F/(k\log_2(F/k))$ against flagged-set density $k/F$, over the sparse-to-moderate range this digest architecture targets. The advantage is steep where flags are sparse ($k\ll F$) and falls toward parity as density grows; the plotted idealization never crosses parity, but the deployed algorithm's constant overhead does invert the ordering at higher density still, exactly as measured at $(F,k)=(100,90)$ in \S\ref{sec:zoom} ($199$ adaptive opens against $100$ flat) --- both halves of the claim are load-bearing.}
\label{fig:r4reg}
\end{figure}
